\chapter{Parties prenantes et utilisateurs ciblés}

\section{Parties prenantes}

Le tableau ci-dessous présente les principales parties prenantes du projet.

\begin{table}[h]
\centering
\begin{tabularx}{\textwidth}{>{\bfseries}l X}
\toprule
\textbf{Partie prenante} & \textbf{Rôle / Intérêt} \\
\midrule
Étudiant acheteur & Recherche des produits, consulte les annonces et souhaite acquérir des articles à moindre coût. \\
Étudiant vendeur & Publie des annonces, modifie ses produits et souhaite valoriser des objets dont il n'a plus besoin. \\
Administrateur & Gère les utilisateurs et modère les annonces signalées. \\
Équipe de développement & Analyse, conçoit, développe et teste la plateforme. \\
Enseignante & Encadre et évalue le projet dans le cadre du module. \\
Communauté étudiante & Bénéficie de la réutilisation des produits et d'un marché de seconde main plus accessible. \\
\bottomrule
\end{tabularx}
\caption{Parties prenantes du projet}
\end{table}

\section{Acteurs du système}

Les acteurs sont les utilisateurs qui interagissent directement avec le système.

\begin{table}[h]
\centering
\begin{tabularx}{\textwidth}{>{\bfseries}l X}
\toprule
\textbf{Acteur} & \textbf{Description} \\
\midrule
Étudiant acheteur & Recherche des produits, consulte les annonces et contacte les vendeurs. \\
Étudiant vendeur & Publie des annonces, modifie ses propres annonces et gère leur disponibilité. \\
Administrateur & Gère les utilisateurs et modère les annonces signalées. \\
\bottomrule
\end{tabularx}
\caption{Acteurs du système}
\end{table}

\textit{Remarque :} Une partie prenante est une personne ou une organisation concernée par le projet, tandis qu'un acteur interagit directement avec le système.
